In this section we introduce girth for closed curves in a handlebody.  We define girth and related concepts and prove some preliminary lemmas.

We use the language of ends.  See~\cite{E} or~\cite{St} for an algebraic introduction to the theory of ends.  See~\cite{GM} for a topological definition.  Let $X$ be a topological space with more than one end.  Define an \emph{interface} for $X$ to be a partition of the ends of $X$ into three subsets labeled black, white, and gray.  We require the black and white subsets of an interface to be nonempty.

A \emph{proper curve} in $X$ is a proper map of the real line to $X$.  Let $\gamma$ be a proper curve in $X$.  The proper curve $\gamma$ determines at most two ends of $X$.  These ends remain invariant under proper homotopies of $\gamma$ (see~\cite[p. 658]{GM}).  Let $\alpha$ be an interface for $X$.  We say $\gamma$ \emph{crosses} $\alpha$ if one end of $\gamma$ in $X$ is black and one is white, as determined by $\alpha$.

Let $H$ be a handlebody.  Let $D$ be an embedded essential disc in $H$.  The disc $D$ defines an interface for the universal cover of $H$, up to deck transformations, as follows.  Let $\hh$ denote the universal cover of $H$.  Consider a lift of $D$ to $\hh$.  This lift separates the ends of $\hh$ into two subsets we call black and white.  Define the gray subset to be empty.  This defines an interface $\alpha$ for $\hh$.  It is well-defined up to deck transformations.  We call $\alpha$ a \emph{lift} of $D$ to an interface for $\hh$.

Let $\gamma$ be a closed curve in $H$.  Let $D$ be an essential disc in $H$.  Define the homotopic intersection number of $\gamma$ and $D$ as follows.  Let $\alpha$ be a lift of $D$ to an interface for $\hh$.  The \emph{homotopic intersection number} of $\gamma$ and $D$ is the number of lifts of $\gamma$ that cross $\alpha$ in the universal cover of $H$.  We denote this quantity $\gamma\wedge D$.  By equivariance, this notion does not depend on the choice of lift $\alpha$.

Homotopic intersection number is invariant under homotopies of $\gamma$.  This is because homotopies of $\gamma$ lift to proper homotopies in the universal cover.  By Lemma~\ref{wedge} below, the homotopic intersection number of $\gamma$ and $D$ equals the minimum geometric intersection number of a curve homotopic to $\gamma$ and $D$.

The \emph{girth} of $\gamma$ in $H$ is the minimum of $\gamma\wedge D$ over all embedded essential discs $D$ in $H$.  We sometimes write this as $\girth(\gamma)$.  An essential disc $D$ is said to \emph{realize} the girth, or be a \emph{girth-realizing} disc, if $\gamma\wedge D=\girth(\gamma)$.  By Lemma~\ref{wedge} below, the girth of $\gamma$ equals the minimum geometric intersection number of an essential disc $D$ in $H$ and a curve homotopic to $\gamma$.

Let $H$ be a handlebody.  Let $\dd$ be a disjoint collection of essential discs in $H$.  Define the homotopic intersection number of $\gamma$ and $\dd$ to be the sum of $\gamma\hi D$ over all discs $D$ in $\dd$.   We denote this $\gamma\hi\dd$.

\begin{lemma}
\label{wedge}
Let $H$ be a handlebody.  Let $\gamma$ be a closed curve in $H$.  Let $\dd$ be a disjoint collection of essential discs in $H$.  The homotopic intersection $\gamma\wedge\dd$ equals the minimum geometric intersection number of $\dd$ and a curve homotopic to $\gamma$.
\end{lemma}

We prove Lemma~\ref{wedge} at the end of this section.  An immediate corollary of Lemma~\ref{wedge} is that if a closed curve $\gamma$ minimizes geometric intersection number with a disc collection $\dd$, then it also minimizes geometric intersection number with any subcollection of $\dd$:

\begin{corollary}
\label{subcollection}
Let $H$ be a handlebody.  Let $\gamma$ be a closed curve in $H$.  Let $\dd$ be a disjoint collection of essential discs in $H$.  Let $\dd_{1}$ be a subcollection of $\dd$.  If $\gamma$ minimizes geometric intersection with $\dd$ over its homotopy class, then $\gamma$ also minimizes geometric intersection with $\dd_{1}$ over its homotopy class.
\end{corollary}

\begin{proof}
Assume $\gamma$ does not minimize intersection with $\dd_{1}$.  It suffices to show that $\gamma$ does not minimize intersection with $\dd$.

Let $\dd_{2}$ denote the collection $\dd-\dd_{1}$.  Let $\gamma\hi\dd_{1}$ equal $n_{1}$, and let $\gamma\hi\dd_{2}$ equal $n_{2}$.  Since $\gamma$ does not minimize intersection with $\dd_{1}$, we have by Lemma~\ref{wedge} that $|\gamma\cap\dd_{1}|>n_{1}$.  Since $\gamma\hi\dd_{2}=n_{2}$, we also have $|\gamma\cap\dd_{2}|\geq n_{2}$, again by Lemma~\ref{wedge}.  Thus $|\gamma\cap\dd|>n_{1}+n_{2}$.

By the definition of homotopic intersection number for a collection of discs, the quantity $\gamma\hi\dd$ equals $n_{1}+n_{2}$.  By Lemma~\ref{wedge}, the minimum geometric intersection number between $\gamma$ and $\dd$ equals $n_{1}+n_{2}$.  So $\gamma$ does not minimize intersection with $\dd$.
\end{proof}

The language of trees will be useful throughout this paper and for the proof of Lemma~\ref{wedge}.  We assume familiarity with notions like paths, geodesics, and backtrackings in trees.  For instance, a geodesic in a tree is a path in the tree with no backtrackings.  See~\cite[Ch. I.2]{Se} for definitions of these terms and an introduction to the theory of trees.

Let $H$ be a handlebody.  Let $\dd$ be a disjoint collection of essential discs in $H$.  We associate an infinite tree $\Gamma$ to $\dd$ in the following way.  Consider all lifts of $\dd$ to the universal cover $\hh$ of the handlebody $H$.  Take the connected components of $\hh-\dd$ to be the vertices of $\Gamma$.  Take the lifts of $\dd$ to $\hh$ to be the edges of $\Gamma$.  Each edge connects two vertices, according to which components of $\hh-\dd$ the lift borders.  We call $\Gamma$ the \emph{tree associated to $\dd$}.

If the discs in $\dd$ are oriented, the graph $\Gamma$ is a directed graph.  If the connected components of $H-\dd$ are simply-connected, then $\Gamma$ is locally finite.  Otherwise, $\Gamma$ is locally infinite.  Here is an important special case.

Let $H$ be a handlebody.  A \emph{disc system} for $H$ is a collection $\dd$ of essential discs in $H$ such that $H-\dd$ is a ball.  If the discs are oriented, we sometime call it an oriented disc system.  If $g$ is the genus of $H$, any disc system has $g$ elements.  An oriented disc system for $H$ naturally gives rise to a collection of free generators for $\pi_{1}(H)$.  In this case, the tree associated to $\dd$ is also the Cayley graph of $\pi_{1}(H)$ with respect to the corresponding generators.  We will sometimes apply Lemma~\ref{wedge} to disc systems.

Let $\gamma$ be a closed curve in $H$ that is transverse to $\dd$.  In this case, we can lift $\gamma$ to a path in $\Gamma$ as follows.  Lift $\gamma$ to $\hh$.  Then map $\gamma$ to $\Gamma$ according to which discs the lift of $\gamma$ crosses.  The resulting path in $\Gamma$ we call a \emph{lift} of $\gamma$ to $\Gamma$.  When $\gamma$ is disjoint from $\dd$, the curve $\gamma$ lifts to a vertex of $\Gamma$, which we view as a degenerate, or constant path.

By the following lemma, $\gamma$ lifts to a path in $\Gamma$ without backtrackings if and only if the curve $\gamma$ minimizes intersection with $\dd$ over its homotopy class.  A path without backtrackings is either an infinite geodesic path or a vertex of $\Gamma$.

\begin{lemma}
\label{geodesic}
Let $H$ be a handlebody.  Let $\dd$ be a disjoint collection of essential discs in $H$.  Let $\Gamma$ be the tree associated to $\dd$.  Let $\gamma$ be a closed curve in $H$ that is transverse to $\dd$.  Lift $\gamma$ to $\Gamma$.  The curve $\gamma$ in $H$ minimizes intersection with $\dd$ in its homotopy class if and only if $\gamma$ in $\Gamma$ has no backtracking.
\end{lemma}

\begin{proof}
Assume first that $\gamma$ in $\Gamma$ has a backtracking.  We will show that $\gamma$ does not minimize intersection with $\dd$.  Choose a backtracking of $\gamma$ in $\Gamma$.  The backtracking corresponds to a component $\kappa$ of $\gamma-D$ in $H$ with the following two properties.  The closure $\overline{\kappa}$ of the arc $\kappa$ intersects $\dd$ only at its endpoints, and the arc $\kappa$ can be homotoped $\rel \bd \kappa$ into a disc in $\dd$.  Homotope $\kappa$ to be disjoint from $D$ in $H$, as shown in Figure~\ref{homotopy 1}.  This reduces the intersection of $\gamma$ and $\dd$ by two, as desired.

Assume now that $\gamma$ lifted to $\Gamma$ has no backtracking.  Then $\gamma$ lifts to either a vertex of $\Gamma$ or an infinite geodesic path in $\Gamma$.  In the former case, $\gamma$ intersects $\dd$ zero times, so it minimizes intersection with $\dd$.  Assume then that $\gamma$ lifts to an infinite geodesic path.  By the first paragraph, it suffices to show that this infinite geodesic is uniquely determined by the homotopy class of $\gamma$.

The infinite path $\gamma$ in $\Gamma$ defines two distinct ends of $\hh$.  Thus $\gamma$ defines two ends of $\Gamma$.   These two ends remain the same under homotopies of $\gamma$ in $H$.  In a tree, an infinite geodesic path is uniquely determined by its ends.  Since $\Gamma$ is a tree, the infinite geodesic determined by $\gamma$ is unique.
\end{proof}

\begin{figure}
\includegraphics[height=2in]{figures/homotopy.eps}
\caption{The gray disc $D$ is part of a disjoint collection of discs $\dd$ in $H$.  Let $\Gamma$ be the tree associated to $\dd$.  If a lift of $\gamma$ to $\Gamma$ has a backtracking (shown above), then the intersection number of $\gamma$ with $\dd$ can be decreased by two in $H$.}
\label{homotopy 1}
\end{figure}
 
We now prove Lemma~\ref{wedge}.

\begin{proof}[Proof of Lemma~\ref{wedge}]
Let $\gamma$ be any curve homotopic to the original $\gamma$.  Then $\gamma\wedge \dd$ is unchanged.  Consider any disc $D$ in $\dd$.  Lift $D$ to the universal cover $\hh$.  The disc $D$ defines an interface $\alpha$ for $\hh$.  If a lift of $\gamma$ crosses $\alpha$, then that lift must geometrically intersect the lift of $D$ at least once.  Thus $\gamma\wedge D \leq |\gamma\cap D|$.  Since this is true for all $D$, we have $\gamma\wedge\dd\leq|\gamma\cap \dd|$.

Now assume $\gamma\wedge \dd<|\gamma\cap\dd|$.  Then we have $\gamma\wedge D<|\gamma\cap D|$ for some $D$ in $\dd$.  It suffices to construct a homotopy of $\gamma$ that reduces the geometric intersection number of $\gamma$ and $\dd$.  Since $\gamma \wedge D<|\gamma\cap D|$, some lift of $\gamma$ must intersect $D$ at least twice.  In particular, $\gamma$ does not lift to a geodesic in $\Gamma$.  By Lemma~\ref{geodesic}, we can decrease the geometric intersection number of $\gamma$ and $\dd$.
\end{proof}